% GENERATED FILE. Edit canonical lessons, then run npm run generate:books.
% canonical-source-hash: fnv1a64:1e856eb64f618895
% canonical-lessons: KA-C65-request, KA-C65-permission, KA-C65-hesitancy, KA-C65-trust, KA-C65-manners

\chapter{More Good Manners}
\label{ch:ka-manners}
\hlchaptermodality{65}

\begin{chapteropening}
\textbf{By the end of this chapter:} \emph{I can name what Kannada courtesy is made of, from a request to good manners.}

Complete KA-C65-manners and say all five words in order.
\end{chapteropening}

\section[vinanti]{\kn{ವಿನಂತಿ} (vinanti) — a request}

\label{lesson:KA-C65-request}



\begin{quote}
Before the new run: what did \emph{bēga} mean, and what did \emph{hāgādare} mean?
\end{quote}

\subsection*{You'll want to know: \kn{ವಿನಂತಿ}}
\textbf{\kn{ವಿನಂತಿ}} (\emph{vinanti}) — \textquotedblleft{}a request\textquotedblright{}.

\kn{ವಿನಂತಿ} is the naming-word for the thing \kn{ದಯವಿಟ್ಟು} has been doing for you since the please-chapter. \kn{ದಯವಿಟ್ಟು} is how you make one; \kn{ವಿನಂತಿ} is what you have made.

Sanskrit, and its shape says \emph{a bending} — the asker lowers himself a little, and that is exactly the difference between a \kn{ವಿನಂತಿ} and an instruction.

\begin{grammarlens}[title={What you've built}]
The first of five courtesy words.
\end{grammarlens}

\subsection*{Your turn}
Say these aloud:

\begin{itemize}
\raggedright
  \item \emph{vinanti}
  \item it once more, slowly
  \item \emph{dayaviṭṭu}, then \emph{vinanti}, and say which one goes inside a sentence
\end{itemize}

\begin{tcolorbox}[breakable,colback=teal!4,colframe=teal!35!black,title={Before you move on}]
What does \kn{ವಿನಂತಿ} mean? (\textquotedblleft{}a request\textquotedblright{}.) Say it once more.
\end{tcolorbox}
\section[anumati]{\kn{ಅನುಮತಿ} (anumati) — permission}

\label{lesson:KA-C65-permission}



\begin{quote}
Before the new one: what did \emph{vinanti} mean?
\end{quote}

\subsection*{You'll want to know: \kn{ಅನುಮತಿ}}
\textbf{\kn{ಅನುಮತಿ}} (\emph{anumati}) — \textquotedblleft{}permission\textquotedblright{}.

\kn{ಅನುಮತಿ} is Sanskrit for a mind that goes along with you — agreement given rather than merely not refused.

It belongs with \kn{ಹೊರಡು} from the leave-taking chapter. In a Kannada house you do not simply set out; you ask \kn{ಅನುಮತಿ} of whoever is host, and only then do you say \kn{ಹೋಗಿ} \kn{ಬರುತ್ತೇನೆ}. The farewell you learned early has this word standing in front of it.

\begin{grammarlens}[title={What you've built}]
Two.
\end{grammarlens}

\subsection*{Your turn}
Say these aloud:

\begin{itemize}
\raggedright
  \item \emph{anumati}
  \item it once more, slowly
  \item \emph{anumati}, then \emph{horaḍu}, and say which one comes first
\end{itemize}

\begin{tcolorbox}[breakable,colback=teal!4,colframe=teal!35!black,title={Before you move on}]
What does \kn{ಅನುಮತಿ} mean? (\textquotedblleft{}permission\textquotedblright{}.) Say it once more.
\end{tcolorbox}
\section[saṅkōca]{\kn{ಸಂಕೋಚ} (saṅkōca) — hesitancy, holding back}

\label{lesson:KA-C65-hesitancy}



\begin{quote}
Before the new one: what did \emph{anumati} mean?
\end{quote}

\subsection*{You'll want to know: \kn{ಸಂಕೋಚ}}
\textbf{\kn{ಸಂಕೋಚ}} (\emph{saṅkōca}) — \textquotedblleft{}hesitancy, holding back\textquotedblright{}.

\kn{ಸಂಕೋಚ} is Sanskrit for a drawing-in, a shrinking — and in Kannada it names the polite reluctance that stops you accepting what you actually want.

It is the one word in this chapter you will hear said to you rather than by you. A host feeding an \emph{atithi} says \kn{ಸಂಕೋಚ} \kn{ಬೇಡ}, using the refusal-word from the replies chapter to refuse your refusal. Hearing it is the signal to eat.

\begin{grammarlens}[title={What you've built}]
Three.
\end{grammarlens}

\subsection*{Your turn}
Say these aloud:

\begin{itemize}
\raggedright
  \item \emph{saṅkōca}
  \item it once more, slowly
  \item \emph{saṅkōca}, then \emph{bēḍa}, and say what a host tells a guest
\end{itemize}

\begin{tcolorbox}[breakable,colback=teal!4,colframe=teal!35!black,title={Before you move on}]
What does \kn{ಸಂಕೋಚ} mean? (\textquotedblleft{}hesitancy, holding back\textquotedblright{}.) Say it once more.
\end{tcolorbox}
\section[nambike]{\kn{ನಂಬಿಕೆ} (nambike) — trust}

\label{lesson:KA-C65-trust}



\begin{quote}
Before the new one: what did \emph{saṅkōca} mean?
\end{quote}

\subsection*{You'll want to know: \kn{ನಂಬಿಕೆ}}
\textbf{\kn{ನಂಬಿಕೆ}} (\emph{nambike}) — \textquotedblleft{}trust\textquotedblright{}.

\kn{ನಂಬಿಕೆ} is native, made out of the everyday word for believing, and it covers both trust in a person and belief in a thing being so.

Set it beside \kn{ಗೌರವ} from the courtesy chapter. \kn{ಗೌರವ} can be given to a stranger at once, because it is owed to a position. \kn{ನಂಬಿಕೆ} has to be collected slowly, and no amount of the first turns into the second.

\begin{grammarlens}[title={What you've built}]
Four.
\end{grammarlens}

\subsection*{Your turn}
Say these aloud:

\begin{itemize}
\raggedright
  \item \emph{nambike}
  \item it once more, slowly
  \item \emph{gaurava}, then \emph{nambike}, and say which one is earned slowly
\end{itemize}

\begin{tcolorbox}[breakable,colback=teal!4,colframe=teal!35!black,title={Before you move on}]
What does \kn{ನಂಬಿಕೆ} mean? (\textquotedblleft{}trust\textquotedblright{}.) Say it once more.
\end{tcolorbox}
\section[saujanya]{\kn{ಸೌಜನ್ಯ} (saujanya) — courtesy, good manners}

\label{lesson:KA-C65-manners}



\begin{quote}
Before the new one: what did \emph{nambike} mean?
\end{quote}

\subsection*{You'll want to know: \kn{ಸೌಜನ್ಯ}}
\textbf{\kn{ಸೌಜನ್ಯ}} (\emph{saujanya}) — \textquotedblleft{}courtesy, good manners\textquotedblright{}.

\kn{ಸೌಜನ್ಯ} is Sanskrit, and underneath it is \emph{a good person} — courtesy named as a quality of the one behaving rather than as a list of rules being followed. That is a real difference, and Kannada means it.

It is the word for everything this book has been quietly teaching alongside the vocabulary: \kn{ಧನ್ಯವಾದ}, \kn{ದಯವಿಟ್ಟು}, \kn{ಕ್ಷಮಿಸಿ}, \kn{ಗೌರವ}, \kn{ವಂದನೆ}, and now the four words in front of this one. \kn{ಸೌಜನ್ಯ} is the name of the whole habit.

\begin{grammarlens}[title={What you've built}]
Five: \kn{ವಿನಂತಿ}, \kn{ಅನುಮತಿ}, \kn{ಸಂಕೋಚ}, \kn{ನಂಬಿಕೆ}, \kn{ಸೌಜನ್ಯ}. Enough to name what you have been doing since the thank-you word.
\end{grammarlens}

\subsection*{Your turn}
Say these aloud:

\begin{itemize}
\raggedright
  \item \emph{saujanya}
  \item it once more, slowly
  \item all five in order, then \emph{dhanyavāda} to close
\end{itemize}

\begin{tcolorbox}[breakable,colback=teal!4,colframe=teal!35!black,title={Before you move on}]
What does \kn{ಸೌಜನ್ಯ} mean? (\textquotedblleft{}courtesy, good manners\textquotedblright{}.) Say it once more.
\end{tcolorbox}
